\section{Data and design}\label{sec:data}

\paragraph{Patents and lapses.} The data are PATSTAT Global \citep[autumn 2025 edition;][]{PATSTAT}. The sample is
granted invention patents filed in 2000--2005 at ten national offices, followed for up to twenty years (to 2019 in Korea and 2022 in Russia, where the
lapse records end). A patent
lapses if the INPADOC legal-event table records its first cessation within 19.5 years of filing; a later cessation
marks expiry at full term. I chose the offices by scanning every authority in the legal-event table for the share of
grants with a recorded cessation and the continuity of records over time (Supplemental Appendix~S1). The
emerging-market offices are those of China, Korea, Russia, Taiwan and the Czech Republic; the advanced-economy
offices are those of the United States, Japan, Germany, the United Kingdom and France. Between them, the five
emerging-market offices granted 84\% of the invention patents filed in 2000--2005 at the national offices of
emerging economies in PATSTAT, and the five advanced-economy offices 88\% of those filed at the national offices of
advanced economies.\footnote{Advanced economies as listed by the IMF in 2005, without Korea and Taiwan; regional offices
such as the European Patent Office are left out.}

\paragraph{Owners.} An owner is domestic if the first applicant is from the office's country. Missing applicant
countries (common in China and Japan) are imputed from the patent family and the office of first filing. Owners are
grouped by sector and size. Size is the number of distinct inventions (DOCDB families) the harmonised applicant filed
worldwide in 2000--2005, the cohort's own filing window, so that it is measured before the outcomes: small firms have
1--9 inventions, mid-sized firms 10--99 and large firms 100 or more; universities and public research organisations
form a fourth group. Individual inventors are excluded: their patents behave differently, and in China they were the
main target of local filing subsidies in this period. Patents granted more than ten years after filing are dropped.
Table~\ref{tab:sample} describes the sample: patents by office and owner in panel A, owners in panel B. It
covers 194{,}518 small firms, 28{,}657 mid-sized firms, 4{,}487 large firms and 17{,}151 universities and public
research organisations, each owner counted once whether it is domestic or foreign.

\begin{table}[t]
\centering
\begin{threeparttable}
\caption{The sample: granted patents filed 2000--2005, firms and universities}\label{tab:sample}
\small
% Source: em-patent-value output/sample_counts.csv (29_sample_tables.py)
\setlength{\tabcolsep}{4pt}
\begin{tabular}{lrrrrrrrrr}
\toprule
 & & \multicolumn{4}{c}{Domestic owners} & \multicolumn{4}{c}{Foreign owners} \\
\cmidrule(lr){3-6}\cmidrule(lr){7-10}
 & & \multicolumn{3}{c}{Firms} & & \multicolumn{3}{c}{Firms} & \\
\cmidrule(lr){3-5}\cmidrule(lr){7-9}
Office & Patents & Small & Mid & Large & Univ. & Small & Mid & Large & Univ. \\
\midrule
\multicolumn{10}{l}{\emph{A. Patents, by office}} \\
China          &   363{,}574 & 26{,}155 & 12{,}411 &  27{,}983 & 39{,}836 & 41{,}845 & 40{,}870 & 143{,}289 &  3{,}982 \\
Korea          &   371{,}332 & 33{,}505 & 21{,}556 & 157{,}473 & 19{,}530 &  8{,}283 & 16{,}088 &  50{,}457 &  1{,}248 \\
Russia         &   129{,}038 & 10{,}130 &  3{,}397 &       841 & 15{,}738 &  8{,}202 &  6{,}144 &   4{,}993 &      467 \\
Taiwan         &   150{,}743 &  6{,}038 &  8{,}190 &  28{,}420 &  6{,}665 &  7{,}510 & 15{,}834 &  76{,}313 &      916 \\
Czech Rep.     &     4{,}061 &      584 &      325 &        -- &       95 &      571 &      826 &   1{,}249 &       24 \\
\addlinespace
United States  &   980{,}998 & 77{,}598 & 87{,}443 & 274{,}701 & 28{,}373 & 47{,}953 & 66{,}673 & 338{,}738 & 17{,}002 \\
Japan          &   775{,}826 & 23{,}780 & 55{,}011 & 539{,}643 & 15{,}540 &  5{,}561 &  9{,}718 &  27{,}163 &  1{,}141 \\
Germany        &    81{,}937 & 11{,}218 & 15{,}326 &  33{,}302 &  3{,}587 &  1{,}669 &  3{,}104 &  13{,}403 &      210 \\
United Kingdom &    41{,}383 &  6{,}095 &  4{,}085 &   1{,}730 &      445 &  3{,}219 &  5{,}028 &  19{,}964 &      591 \\
France         &    62{,}937 & 11{,}228 & 10{,}448 &  21{,}569 &  3{,}753 &  1{,}465 &  1{,}957 &   8{,}368 &      162 \\
\midrule
Total          & 2{,}961{,}829 & 206{,}331 & 218{,}192 & 1{,}085{,}662 & 133{,}562 & 126{,}278 & 166{,}242 & 683{,}937 & 25{,}743 \\
\midrule
\multicolumn{10}{l}{\emph{B. Owners, by office}} \\
% Source: em-patent-value output/owner_counts.csv (29b_owner_counts.py)
China          &   & 18{,}838 & 1{,}014 & 157 & 6{,}122 & 30{,}420 & 7{,}518 & 2{,}084 & 1{,}738 \\
Korea          &   & 18{,}996 & 1{,}798 & 243 & 732 & 5{,}607 & 3{,}730 & 1{,}592 & 433 \\
Russia         &   & 6{,}121 & 176 & 5 & 3{,}679 & 4{,}735 & 407 & 93 & 289 \\
Taiwan         &   & 3{,}513 & 563 & 82 & 183 & 5{,}255 & 3{,}295 & 1{,}209 & 270 \\
Czech Rep.     &   & 373 & 19 & -- & 55 & 472 & 465 & 291 & 23 \\
\addlinespace
United States  &   & 42{,}360 & 5{,}853 & 787 & 1{,}369 & 32{,}868 & 8{,}435 & 1{,}448 & 2{,}129 \\
Japan          &   & 15{,}610 & 6{,}588 & 1{,}684 & 497 & 3{,}674 & 1{,}284 & 510 & 282 \\
Germany        &   & 7{,}270 & 1{,}993 & 299 & 341 & 1{,}246 & 830 & 545 & 79 \\
United Kingdom &   & 3{,}916 & 531 & 44 & 145 & 2{,}399 & 1{,}451 & 767 & 210 \\
France         &   & 6{,}292 & 739 & 99 & 296 & 1{,}071 & 634 & 474 & 60 \\
\midrule
All offices    &   & 123{,}285 & 19{,}258 & 3{,}365 & 13{,}410 & 78{,}990 & 16{,}172 & 3{,}015 & 4{,}510 \\
\bottomrule
\end{tabular}
\begin{tablenotes}\footnotesize
\item \emph{Notes:} Granted invention patents filed 2000--2005 with a grant lag of up to ten years; individual
owners excluded. Patents: all firms, universities and owners of unknown type. Firm size: inventions (DOCDB families)
filed worldwide by the harmonised applicant in 2000--2005; small 1--9, mid-sized 10--99, large 100 or more; no Czech
firm reaches 100. Univ.: universities and public research organisations. Panel B: distinct harmonised applicants with at
least one sample patent in the office. An owner patenting in several offices is counted in each; All offices counts
it once among domestic and once among foreign owners, so the domestic and foreign totals overlap and fall short
of the sums over offices (an owner whose patents name applicants from different countries can be domestic in more
than one office). Source: PATSTAT Global, autumn 2025.
\end{tablenotes}
\end{threeparttable}
\end{table}



\paragraph{Quality.} Observable quality is measured by forward citations to the patent family (logged) and the
WIPO technology field (35 fields), the controls in the main estimates. Robustness checks add citations in search
reports that block later claims (categories X and Y), citations received in the four years after filing, the grant
lag and family size. Domestic patents in emerging-market offices are cited far less than foreign ones,
by a factor of five in China and six in Korea, so quality controls matter.

\paragraph{Design.} The basic comparison is between domestic and foreign owners in the same office, which holds
fees, law and enforcement fixed. Foreign owners are mostly multinationals protecting a foreign market, so domestic
and foreign patents differ everywhere; the question is whether the domestic shortfall is larger in emerging-market
offices. For patent $i$ of owner group $g$ in office $o$, filed in year $t$ in field $f$,
\begin{equation}\label{eq:did}
  y_{i} = \beta\, D_i + \delta\, D_i \times EM_o + \theta_o \ln(1+c_i) + \mu_{ot} + \lambda_{of} + \kappa_{og}
  + \varepsilon_i ,
\end{equation}
where $y_i$ indicates that the patent is alive eight years after grant,\footnote{Survival is measured from
grant because renewal fees are due from grant. Grant lags differ across offices and between domestic and foreign
owners in the same office (Supplemental Appendix Figure~S1), so a fixed age from filing would
compare patents that have faced different numbers of renewal decisions.} $D_i$ marks a domestic owner, $EM_o$ an
emerging-market office and $c_i$ citations. The coefficient of interest, $\delta$, is the extra domestic--foreign gap
in emerging-market offices, in percentage points. Standard errors are clustered by owner.
